%% en/sections/05_npc.tex — NPCs EN (full)

\chapter{Non-Player Characters}
\markboth{Non-Player Characters}{Non-Player Characters}

\lettrine[lines=3]{N}{on-player} characters in \textit{Vespri 1282} are,
where possible, historically attested figures. Each is presented with
a game block and a historiographical note to help GMs portray them
consistently. Historical NPCs are anchors, not cages: amplify, reduce,
or reposition them as needed --- just ask whether the choice is coherent
with what we know of that person.

\secsep

\section{Cross-Faction NPCs}

\subsection{Giovanni da Procida}
\textit{Physician, jurist, diplomat. The architect of the conspiracy — or its legend.}

\statblock{Giovanni da Procida}{Lead Conspirator}{3}{4}{1}{3}

\textbf{Available to:} Conspirators (direct), Barons (emissary),
Clergy (indirect contact)

Da Procida is not in Palermo --- he is the distant reference point
of the conspiracy. He appears through letters, emissaries, or a special
scene if the GM chooses. His presence is most powerful when he is absent:
characters discuss him, interpret his instructions, wonder whether he
trusts them.

\textbf{Character:} Cold, methodical, capable of genuine affection only
for the cause. He knows exactly what the people he uses are worth ---
and uses them anyway.

\textbf{Historical note:} His physical presence in Palermo in 1282 is
disputed. Some sources place him in Sicily; others in Aragon. In play,
this ambiguity is productive.

\filetto

\subsection{Jean de Saint-Rémy}
\textit{The Angevin justiciar of Palermo. He is doing his job. That is the problem.}

\statblock{Jean de Saint-Rémy}{Royal Justiciar}{2}{3}{2}{2}

\textbf{Available to:} All tables (primary or secondary antagonist)

De Saint-Rémy is not a textbook villain. He is a competent official
doing his job in a situation that is escaping his control. He knows
something is moving, but does not know where to look. He is paranoid,
reactive, and willing to use violence --- but prefers prevention to
repression.

\textbf{Special mechanic --- The De Saint-Rémy Clock:} The coordinating
GM keeps a hidden 6-segment clock. Every time a table takes an action
that draws Angevin attention, the clock advances. When full: curfew,
port closure, reinforced patrols, -1 to every faction score.

\filetto

\subsection{Archbishop of Palermo}
{\cinzel\footnotesize\color{grigionote} Guglielmo di Monforte}\\
\textit{Between God, Rome, and his flock. Three masters. None satisfied.}

\statblock{Archbishop of Palermo}{Prelate — Guglielmo di Monforte}{2}{3}{1}{4}

\textbf{Available to:} Clergy (direct), Barons (negotiation),
Conspirators (indirect)

The Archbishop knows more than he shows. He wants his people to fare
well. He does not want to be complicit in mass killing. These two things
may be incompatible.

\secsep

\section{Conspirators NPCs}

\subsection{Aragonese Emissary}
\statblock{The Emissary}{Aragonese Agent — Guillem de Montcada}{2}{3}{2}{2}

Liaison between the Sicilian network and the court of Barcelona.
Carries instructions, gold, and information --- not all of it accurate.
Peter of Aragon has his own timing, which does not necessarily match
the uprising's. The emissary is where this friction becomes concrete.

\textbf{Secret:} Knows Peter may not intervene militarily if the
uprising breaks out at the wrong time. Has no instructions on how
to manage this information.

\subsection{Messina Contact}
\statblock{Messina Contact}{Merchant-Agent — Enrico di Taormina}{2}{2}{2}{3}

Holds the eastern Sicily network. If the Conspirators keep him
operational, the uprising can extend to Messina --- doubling its
impact. If the network is compromised, he vanishes.

\subsection{Byzantine Agent}
\statblock{Byzantine Agent}{Diplomat-Spy — Alexios Philanthropenos}{1}{4}{1}{3}

Represents the interests of Michael VIII Palaeologus. Finances the
conspiracy but wants guarantees that Charles's expedition will be
blocked --- regardless of what happens to Sicily. His interests and
the Conspirators' overlap without coinciding.

\secsep

\section{Barons NPCs}

\subsection{Alaimo da Lentini}
\textit{The most powerful baron of eastern Sicily. And the most cautious.}

\statblock{Alaimo da Lentini}{Baron}{3}{3}{2}{3}

\textbf{Available to:} Barons (reference figure), Conspirators
(potential ally)

Alaimo is the weight that could shift the balance --- but he does not
move easily. He has seen too many plans fail to trust the first one
that comes along. A convinced Alaimo automatically adds +1 to the
faction score.

\textbf{Historical note:} Alaimo da Lentini is attested as one of the
uprising's leaders in eastern Sicily. He was later implicated in a
counter-conspiracy against Peter of Aragon and executed in 1287.
His real narrative arc is a tragedy.

\subsection{Pro-Angevin Vassal}
\statblock{Pro-Angevin Vassal}{Minor Baron — Ugo di Santapau}{2}{2}{2}{2}

Internal antagonist. Not a declared enemy, but his choices inform
the French --- consciously or not. The Barons must decide what to do
with him: neutralise, convert, or ignore.

\subsection{Rival Family}
\statblock{Rival Family Head}{Baron — Ruggero di Incisa}{2}{2}{3}{2}

An old feud complicates baronial unity. An impossible alliance,
if achieved, is worth +1 to the score.

\secsep

\section{Clergy NPCs}

\subsection{Fra' Matteo}
\statblock{Fra' Matteo di Sant'Agata}{Minor Friar}{3}{2}{2}{3}

Natural informant for the Clergy. Moves through the Albergheria and
Kalsa quarters, knows the families, knows where French soldiers are
in the evening. Not a conspirator --- a person who wants his people
to fare well.

\subsection{Papal Legate}
\statblock{Papal Legate}{Papal Emissary — Gherardo da Parma}{1}{4}{1}{4}

Complex antagonist for the Clergy. Not an enemy of Sicily --- loyal
to Rome, which currently aligns with Charles. If the Clergy convinces
him of what is really happening, he might choose not to intervene.
He will not help --- but he might not obstruct.

\subsection{The Visionary Nun}
\statblock{Sister Eufemia di Corleone}{Nun}{2}{2}{1}{4}

Ambiguous figure. Her visions are taken seriously by the population.
If the Clergy frames them narratively (the visions speak of liberation),
they become a powerful tool. If they lose control, they may trigger
an uncontrolled reaction.

\secsep

\section{People NPCs}

\subsection{The French Soldier}
\textit{He is twenty. He is far from home. He is also afraid.}

\statblock{French Soldier}{Angevin Militiaman — Renaud de Coucy}{2}{1}{3}{1}

Not a monster. A young man following orders in a country he does not
understand. People PCs may choose to hate him, pity him, or use him.
Each choice has different consequences.

\nota[Design note]{
  This NPC is deliberately built to create moral discomfort.
  Violence against him must carry an emotional cost.
}

\subsection{The Quarter Chief}
\statblock{Quarter Chief}{Community Figure — Salamone dell'Albergheria}{3}{2}{2}{3}

Potential People ally. Not a revolutionary --- a practical man who
wants his quarter to survive. If the PCs convince him the uprising
is inevitable, he will organise the quarter. If not, he will take cover.

\subsection{The Mother of the Dead}
{\cinzel\footnotesize\color{grigionote} Addolorata di Seralcadio}\\
\textit{Her son was killed by the French three months ago. No one answered. She has not forgotten.}

\statblock{The Mother of the Dead}{Witness — Addolorata di Seralcadio}{1}{2}{1}{4}

Not an agent --- a living testimony. If the PCs bring her publicly
to tell her story, the next Incite roll is automatically Great effect.
Her pain is real. Using it is a moral choice the PCs must make
consciously.

\filettodoppio
